\documentclass{article}
\usepackage[left=2cm, right=2cm, top=2cm, bottom=2cm]{geometry}
\usepackage{graphicx}
\usepackage{amsmath}
\usepackage{amssymb}
\usepackage{amsthm}
\usepackage{fancyhdr}
\usepackage{verbatim}
\usepackage{listings}
\usepackage{xcolor}
\usepackage{pdfpages}
\usepackage{cancel}
\usepackage{physics}
\usepackage{esint}

\lstset{
    language=C++,
    basicstyle=\ttfamily\footnotesize,
    keywordstyle=\color{blue},
    commentstyle=\color{green},
    stringstyle=\color{red},
    numbers=left,
    numberstyle=\tiny\color{gray},
    frame=single,
    breaklines=true,
    captionpos=b,
}


\title{PHYS 487: Lecture \# 20}
\author{Cliff Sun}

\newtheorem{theorem}{Theorem}[section]
\newtheorem{lemma}[theorem]{Lemma}
\newtheorem{definition}[theorem]{Definition}
\newtheorem{conjecture}[theorem]{Conjecture}
\newtheorem{proposition}[theorem]{Proposition}
\newtheorem{corollary}[theorem]{Corollary}
\newtheorem{one minute paper}[theorem]{One Minute Paper}

\pagestyle{fancy}
\lhead{\textbf{\thepage}\ \ \nouppercase{\rightmark}}
\chead{PHYS 487: Lecture \# 20}
\rhead{Cliff Sun}

\begin{document}

\maketitle

\section*{Lecture Span}
\begin{enumerate}
    \item Introduction to Quantum Information
\end{enumerate}

\section*{Intro to QI}

Today, we will understand the root cause of complexity in quantum information theory. Also, does this "complexity" in quantum information help us with computing? 
Also, how does quantum become "classical", i.e. measurement causing a transition. 

Note, quantum information theory is the intersection of information theory and quantum mechanics. Thus, quantum computation is the intersection of QM predictions and QI. 

Now, consider $N$ qubits with states $\alpha_i\ket{0}_i + \beta_i\ket{1}_i$. In general, a regular computer requires about 64 bits of information to store $\alpha_i, \beta_i$. Therefore, 
a computer needs exponentially large amount of bits.

If $N \sim 100$+, then it is completely classically untractable. (need more bits than protons in the visible universe) 

\section*{Quantum Circuits}

Consider a initial state $\ket{0}$. Then imagine that we apply a time evolution operator $u$ on $\ket{0}$ and measure $\ket{0}$. The time evolution operator can be decomposed into a series of operations done by 
$\sigma_i$ (Pauli Matrices). These are also "gates". We can consider quantum circuits with multiple qubits.

Consider $\ket{0}_a$ and perform a unitary $u_a$. Similarly for $\ket{0}_b$ and $u_b$. Then, we can introduce multi-qubit gates. Such gates can cause entangling, such as the CNOT gate (controlled not). (if the control qubit is $\ket{1}$, then invert the state of the other qubit). 

\subsection*{Example with superposition input}

Let $a: \ket{+}$ and $b: \ket{0}$. Here, $$\ket{+} = \frac{1}{\sqrt{2}}\left( \ket{0} + \ket{1} \right)$$

Then act a CNOT gate on both of the qubits and measure both qubits. Here, 

\begin{align}
    \ket{\psi}_{ab} &= \frac{1}{\sqrt{2}}\left( \ket{0}_a + \ket{1}_a\right)\ket{0}_b \\
    &= \frac{1}{\sqrt{2}}\left( \ket{0}\ket{0} + \ket{1}\ket{0} \right)
\end{align}

Then, acting the CNOT gate on this system yields

\begin{equation}
    \frac{1}{\sqrt{2}}\left( \ket{0}\ket{0} + \ket{1}\ket{1} \right)
\end{equation}

This is an entangled state (bell state). Note, no matter what basis we measure this state in, we will always observe correlations. 

\section*{Deutsch's Algorithm}

Consider a one bit function $f(x)$ where $x$ is a classical bit of which we input into $f$. There are 4 total possible functions that we can define. 
Then, Deutsch's Algorithm solves if $f$ is a constant function. 

Classically, we need to call this function twice in order to determine if $f(0) = f(1)$. But quantumly, we only have to call the function once. Consider a two qubit system $\ket{x}$ and $\ket{y}$ where y is the result. 

Then act $u_f$ on $x$ and $y$ with $x$ is the same and $y \rightarrow y \oplus f(x)$. 

\end{document}